\documentclass[10pt]{beamer}

\usepackage{appendixnumberbeamer}
\usepackage{booktabs}
\usepackage[scale=2]{ccicons}
\usepackage{pgfplots}
\usepackage{xspace}
\usepackage{bookmark}
\usepackage{amssymb}
\usepackage{mathtools}
\usepackage[normalem]{ulem}
\usepackage[T1]{fontenc}
\usepackage[sfdefault,book]{FiraSans}
\usepackage{FiraMono}
\usepackage{hyperref}
\usetheme{metropolis}

\usepgfplotslibrary{dateplot}

\lstset{
  basicstyle=\ttfamily,
  escapeinside={\%*}{*)},
  mathescape=true,
  literate=
  {á}{{\'a}}1 {é}{{\'e}}1 {í}{{\'i}}1 {ó}{{\'o}}1 {ú}{{\'u}}1
  {Á}{{\'A}}1 {É}{{\'E}}1 {Í}{{\'I}}1 {Ó}{{\'O}}1 {Ú}{{\'U}}1
  {à}{{\`a}}1 {è}{{\`e}}1 {ì}{{\`i}}1 {ò}{{\`o}}1 {ù}{{\`u}}1
  {À}{{\`A}}1 {È}{{\'E}}1 {Ì}{{\`I}}1 {Ò}{{\`O}}1 {Ù}{{\`U}}1
  {ä}{{\"a}}1 {ë}{{\"e}}1 {ï}{{\"i}}1 {ö}{{\"o}}1 {ü}{{\"u}}1
  {Ä}{{\"A}}1 {Ë}{{\"E}}1 {Ï}{{\"I}}1 {Ö}{{\"O}}1 {Ü}{{\"U}}1
  {â}{{\^a}}1 {ê}{{\^e}}1 {î}{{\^i}}1 {ô}{{\^o}}1 {û}{{\^u}}1
  {Â}{{\^A}}1 {Ê}{{\^E}}1 {Î}{{\^I}}1 {Ô}{{\^O}}1 {Û}{{\^U}}1
  {Ã}{{\~A}}1 {ã}{{\~a}}1 {Õ}{{\~O}}1 {õ}{{\~o}}1
  {œ}{{\oe}}1 {Œ}{{\OE}}1 {æ}{{\ae}}1 {Æ}{{\AE}}1 {ß}{{\ss}}1
  {ű}{{\H{u}}}1 {Ű}{{\H{U}}}1 {ő}{{\H{o}}}1 {Ő}{{\H{O}}}1
  {ç}{{\c c}}1 {Ç}{{\c C}}1 {ø}{{\o}}1 {å}{{\r a}}1 {Å}{{\r A}}1
  {€}{{\euro}}1 {£}{{\pounds}}1 {«}{{\guillemotleft}}1
  {»}{{\guillemotright}}1 {ñ}{{\~n}}1 {Ñ}{{\~N}}1 {¿}{{?`}}1
}

\title{Desenvolvimento de Software para a Web II}
\author{Thiago Cavalcante  -- thiago.cavalcante@penedo.ufal.br}
\institute{Universidade Federal de Alagoas -- UFAL \\ Campus Arapiraca \\ Unidade de Ensino de Penedo}
\titlegraphic{\hfill\includegraphics[height=1.5cm]{images/brasao-ufal.eps}}


\subtitle{Aula 1}
\date{24 de outubro de 2019}

\begin{document}

\maketitle

\begin{frame}[fragile]{Informações Gerais da Disciplina}
    \begin{center}
        \begin{tabular}{@{}ll@{}}
            \toprule
            \textbf{Nome} & Desenvolvimento de Software para a Web II \\
            \textbf{Código} & SISB031 \\
            \textbf{Semestre} & 7º \\
            \textbf{Carga Horária} & 60h \\
            \textbf{PPC} & 03/2018 \\
            \textbf{Turma} & 2019.2 \\
            \textbf{Horário} & Quarta, 19:50 -- 22:30 \\
            \textbf{Sala} & 2 \\
            \textbf{Lista de discussão} & Google Groups: \href{https://groups.google.com/forum/#!forum/sisb031_20192}{\alert{sisb031\_20192}} \\
            \textbf{Repositório} & GitHub: \href{https://github.com/theagoliveira/sisb031_20192}{\alert{theagoliveira/sisb031\_20192}} \\
            \bottomrule
        \end{tabular}
    \end{center}
\end{frame}

\begin{frame}{Cronograma}
    \begin{center}
        \begin{tabular}{@{}r@{/}r@{/}rll@{}}
            \toprule
            \multicolumn{4}{@{}l}{\textbf{Datas importantes}} \\
            \midrule
            21 & 11 & 2019 & \textit{Possível} dia da \alert{AB1} \\
            23 & 11 & 2019 & Prazo final para digitação da \alert{AB1} \\
            13 & 02 & 2020 & \textit{Possível} dia da \alert{AB2} \\
            17 & 02 & 2020 & Prazo final para digitação da \alert{AB2} \\
            17--22 & 02 & 2020 & Período de \alert{reavaliação} \\
            27--29 & 02 & 2020 & Período de \alert{provas finais} \\
            \bottomrule
        \end{tabular}
    \end{center}
\end{frame}

\begin{frame}{Recapitulando: Desenv. de Software para a Web I}
    \begin{itemize}
        \item Introdução ao ambiente WEB
        \item Projetos de sistemas para web: modelo cliente-servidor, padrão MVC, arquitetura em camadas, protocolo HTTP
        \item Introdução e fundamentação ao HTML: textos, paragráfos, containers, hiperlinks, tabelas, imagens, formulários e menus
        \item Introdução às folhas de estilo em cascata (CSS)
        \item Conceitos de scripts do lado do cliente com Javascript e suas aplicações em páginas web
        \item Introdução a Web dinâmica
        \item Desenv. de aplicações web usando \textit{Server Side Scripts} (PHP)
        \item Gerenciamento de estado de aplicações web
        \item Criação de formulários web
        \item Desenvolvimento de aplicações com banco de dados
    \end{itemize}
\end{frame}

\begin{frame}{Ruby on Rails}
    \includegraphics[width=\textwidth]{images/rails_logo.png}
\end{frame}

\setbeamertemplate{frame footer}{https://ideamotive.co/blog/40-best-ruby-on-rails-companies-websites/}
\begin{frame}{Por que Rails?}
    Top 10 empresas que usam Rails em 2019:

    \begin{itemize}
        \item AirBNB
        \item Groupon
        \item GitHub
        \item Couchsurfing
        \item Shopify
        \item Ask.fm
        \item Dribble
        \item Twitter
        \item Etsy
        \item Fab
    \end{itemize}
\end{frame}
\setbeamertemplate{frame footer}{}

\begin{frame}{Bibliografia}
    \huge
    Codecademy: Learn Ruby

    \bigskip

    Michael Hartl: Ruby on Rails Tutorial
\end{frame}

\begin{frame}{Conceitos de Ruby}
    \Large
    \begin{center}
        \includegraphics[width=0.5\textwidth]{images/ruby_logo.png}
    \end{center}

    Link para programar em Ruby no navegador:

    https://repl.it/languages/ruby
\end{frame}

\end{document}
